\Page{4}{\textenglish{EDITOR’S NOTE}}{החנות שלפני החנות}{46}{53}

{\fontsize{13.2}{18.7}\selectfont

המתלה בחנות מציג תוצאה שכבר נבחרה. ליד החולצה הלבנה יש גרסה ירוקה, ומאחוריה מכנסיים שמתאימים בדיוק לגובה שבו היא נגמרת. הכל נראה מובן מאליו. כמה חודשים קודם, הירוק היה אולי אחת מכמה הצעות, למכנסיים היו שתי צלליות אפשריות, ולשום פריט עדיין לא הייתה תמונה שהקהל יכול לגלול אליה. בינתיים מישהי הייתה צריכה לבחור מה ייכנס למגוון, איך להראות אותו ואיפה להציע אותו. התעשייה שמייצרת את הרגע הקטן הזה בחנות משקיעה עבודה רבה בדברים שאף אחת מהקונות לא תראה.

\par}

\Photo{discovery.png}{92mm}{צלליות אופנה בעיר. \AIArt}

{\fontsize{12.7}{18}\selectfont

\StoryColumns{עכשיו גם הבחירות המוקדמות נעשות מול חומר עבודה אחר. אפשר להעמיד שתי צלליות על דמות ממוחשבת לפני שתפרו דוגמה. אפשר לדעת אם צבע מופיע יותר בתמונות, ולהמשיך לבדוק אם העלייה רלוונטית לקהל בחנות. אפשר להחליף בתמונה סביבה שלמה, ועדיין לעבור פרט אחר פרט על הבגד שיימכר. כשהאפשרויות נעשות קלות יותר להצגה, עולה הערך של מי שיודעת להבחין ביניהן: לזהות איזה הבדל שווה ניסוי, ואיזה הבדל נובע רק מהדרך שבה התוכנה הציגה את הדברים. \SourceRef{2}\SourceRef{3}\SourceRef{4}\SourceRef{5}}{אפשר להרגיש את השינוי גם בלי להיכנס לסטודיו. תמונת מוצר יכולה להפוך להדמיה על תמונה אישית. חשבון שנראה כמו יומן של משפיענית עשוי להיות ערוץ שמותג מפיק ומנהל. שתי החוויות יכולות להיות משכנעות גם כשהן דורשות מאיתנו לקרוא את התמונה באופנים שונים. זו הסיבה להתעניין בצד המקצועי שלהן. הוא נותן לקוראת דרך להבין אילו החלטות התקבלו עבורה לפני שהגיעה לחנות, ואיזה חלק מהבחירה עדיין נשאר בידיה. \SourceRef{10}\SourceRef{11}}

\par}
